% 算法题LaTeX模板 - 双栏紧凑排版
% 编译方式: xelatex main.tex (推荐) 或 pdflatex main.tex

\documentclass[10pt, a4paper]{article}

% ============ 基础包 ============
\usepackage[utf8]{inputenc}
\usepackage[T1]{fontenc}
\usepackage{ctex}           % 中文支持
\usepackage{geometry}       % 页面设置
\usepackage{multicol}       % 多栏布局
\usepackage{enumitem}       % 列表设置
\usepackage{listings}       % 代码高亮
\usepackage{xcolor}         % 颜色
\usepackage{titlesec}       % 标题间距
\usepackage{parskip}        % 段落间距
\usepackage{amsmath}        % 数学公式
\usepackage{graphicx}       % 图片

% ============ 极限页面设置 ============
\geometry{
    left=0.5cm,
    right=0.5cm,
    top=0.5cm,
    bottom=0.5cm,
    nohead,    % 无页眉
    nofoot     % 无页脚
}

% ============ 代码样式设置 ============
\definecolor{codebg}{RGB}{245,245,245}
\definecolor{codegreen}{RGB}{0,150,0}
\definecolor{codeblue}{RGB}{0,0,200}
\definecolor{codered}{RGB}{200,0,0}

\lstset{
    basicstyle=\ttfamily\scriptsize,     % 极小字号
    breaklines=true,                     % 自动换行
    breakatwhitespace=false,
    breakindent=0pt,
    frame=single,                        % 单线框
    backgroundcolor=\color{codebg},      % 背景色
    keywordstyle=\color{codeblue}\bfseries,
    commentstyle=\color{codegreen},
    stringstyle=\color{codered},
    showstringspaces=false,
    captionpos=b,
    aboveskip=2pt,                       % 上方间距
    belowskip=2pt,                       % 下方间距
    xleftmargin=2pt,
    xrightmargin=2pt,
    framesep=3pt,
    framerule=0.5pt,
    tabsize=4,
    numbers=none,                        % 不显示行号
}

% ============ C++ 代码设置 ============
\lstdefinestyle{cppstyle}{
    language=C++,
    morekeywords={include, define, ifdef, ifndef, endif, using, namespace, std, vector, string, map, set, unordered_map, unordered_set},
}

% ============ 紧凑标题设置 ============
\titlespacing*{\section}{0pt}{4pt}{2pt}  % 左边距, 上间距, 下间距
\titlespacing*{\subsection}{0pt}{2pt}{1pt}
\titlespacing*{\subsubsection}{0pt}{1pt}{0pt}

% 设置标题字号
\setcounter{secnumdepth}{0}  % 不显示编号
\setlength{\columnsep}{0.3cm}  % 栏间距

% 自定义section格式
\titleformat{\section}
  {\normalfont\bfseries\large}{\thesection}{0pt}{}

\titleformat{\subsection}
  {\normalfont\bfseries\normalsize}{}{0pt}{}

% ============ 全局样式调整 ============
\raggedcolumns  % 栏内容不拉伸对齐，保持自然紧凑
\setlength{\parindent}{0pt}      % 无首行缩进
\setlength{\parskip}{3pt}        % 段落间距
\setitemize{itemsep=1pt,topsep=1pt,leftmargin=1.2em}  % 列表紧凑

\begin{document}

% ============ 标题页 ============
\begin{center}
    \huge\textbf{分治法}
\end{center}

\vspace{0.3cm}
\noindent\rule{\linewidth}{0.4pt}
\vspace{0.3cm}

% ============ 目录页 ============
\begin{multicols*}{2}
% ========== 自动生成的目录内容 ==========
\vspace{8pt}
1. \textbf{寻找凸包（Graham扫描法）}\hfill 2\\
给定平面上的点集Q，使用Graham扫描法求Q的凸包（凸包定义：包含所有点的最小凸多边形，点或在多边形上或在内部），要求展示算法的完整步骤。

\vspace{8pt}
2. \textbf{两个有序数组的中位数}\hfill 2\\
设X[0:n-1]和Y[0:n-1]为两个长度均为n的有序数组，设计O(log n)时间的分治算法，找出X和Y的2n个元素的中位数（偶数个元素取中间两个的平均值，奇数个取最中间元素）。

\vspace{8pt}
3. \textbf{最大子数组问题}\hfill 3\\
给定包含n个整数（有正有负）的数组A，设计O(nlogn)时间的分治算法找出和最大的非空连续子数组，同时设计O(n)时间的Kadane算法。

\vspace{8pt}
4. \textbf{循环移位问题}\hfill 3\\
给定一个按从小到大排序的数组，现将其循环右移若干位（如原数组[1,2,3,4,5]，循环右移2位后为[4,5,1,2,3]），设计算法计算循环右移的位数k（$k \ge 0$），要求时间复杂度尽可能优化。

\vspace{8pt}
5. \textbf{归并排序}\hfill 4\\
给定数组A，使用归并排序算法对其进行升序排序，要求展示分治的"分解$\to$求解$\to$合并"完整过程。

\vspace{8pt}
6. \textbf{折半搜索（查找成功案例）}\hfill 4\\
给定有序数组A，使用折半搜索算法查找目标元素target，要求展示每次查找的low、mid、high值及比较过程，若查找成功返回元素下标，失败返回-1。

\vspace{8pt}
7. \textbf{折半搜索（查找失败案例）}\hfill 5\\
给定有序数组A，使用折半搜索算法查找不存在的目标元素，重点关注查找失败时的边界调整及终止条件判断。

\vspace{8pt}
8. \textbf{大整数乘法（Karatsuba算法）}\hfill 5\\
设X和Y为两个n位大整数（n为偶数），X拆分为A（前n/2位）和B（后n/2位），Y拆分为C（前n/2位）和D（后n/2位），即$X=A \cdot 2^{n/2}+B$，$Y=C \cdot 2^{n/2}+D$。设计优化方案通过减少乘法次数降低时间复杂度。

\vspace{8pt}
9. \textbf{Strassen矩阵乘法}\hfill 6\\
给定两个n$\times$n矩阵A和B，使用Strassen算法（分治法）计算矩阵乘积C=A$\times$B，要求展示7个中间矩阵的定义及复杂度分析。

\columnbreak
\vspace{8pt}
10. \textbf{快速排序}\hfill 6\\
给定数组A，使用快速排序算法对其进行升序排序，要求展示Partition过程及完整的递归排序步骤。

\vspace{8pt}
11. \textbf{第k小元素（选择问题）}\hfill 7\\
给定无序数组A，使用分治法查找该数组的第k小元素，要求展示Partition过程及递归查找步骤。

\vspace{8pt}
12. \textbf{堆排序}\hfill 7\\
给定数组A，使用堆排序算法对其进行升序排序，要求展示最大堆的建立过程、堆顶元素的弹出与堆调整过程。

\vspace{8pt}
13. \textbf{最近点对问题}\hfill 8\\
给定平面上n个点的集合Q，设计分治算法找出距离最近的两个点，要求说明预处理、分解、递归求解、合并的完整流程。

\vspace{8pt}
14. \textbf{天际轮廓问题}\hfill 8\\
给定n座建筑物，每个建筑物用三元组(Li,Ri,Hi)表示（Li为左下x坐标，Ri为右下x坐标，Hi为高度），设计O(nlogn)算法求出这n座建筑物的天际轮廓。

\vspace{8pt}
15. \textbf{两元素和为X的分治算法}\hfill 9\\
给定由n个实数构成的集合S和实数X，判断S中是否存在两个元素的和为X，要求设计分治算法求解。

\vspace{8pt}
16. \textbf{逆序对的分治算法}\hfill 9\\
给定实数序列A，若i<j且A[i]>A[j]，则(i,j)为逆序对。要求用分治方法求序列中逆序对的个数。

\vspace{8pt}
17. \textbf{数组中两元素和大于0的对数量}\hfill 10\\
给定大小为n的数组A，求数组中不同对(i,j)（i<j）的数量，使得A[i]+A[j]>0。

\vspace{8pt}
18. \textbf{分治算法解题的固有方法论}\hfill 11\\
分治算法（Divide and Conquer）的核心思想是"大事化小，小事化了，最后合并"。标准三步走策略：分解(Divide)、解决(Conquer)、合并(Combine)。复杂度分析使用主定理 $T(n) = aT(n/b) + f(n)$。

\end{multicols*}

\vspace{0.5cm}

% ============ 双栏开始 ============
\begin{multicols*}{2}

% ========== 题目内容 ==========
\section*{1. 寻找凸包（Graham扫描法）}

给定平面上的点集Q，使用Graham扫描法求Q的凸包（凸包定义：包含所有点的最小凸多边形，点或在多边形上或在内部），要求展示算法的完整步骤。

\begin{itemize}
    \item \textbf{基准点选择}：找y坐标最小的点p0（y相同时取x最小），保证p0必在凸包上
    \item \textbf{极角排序}：其余点按相对p0的极角从小到大排序，共线点按距离由近及远，用叉积判断避免角度计算
    \item \textbf{栈维护凸性}：从第3个点开始扫描，若新点与栈顶两点不构成左转（叉积$\le 0$）则弹出栈顶，直至构成左转或栈只剩1个点
    \item \textbf{时间复杂度}：找基准点O(n)，极角排序O(nlogn)，扫描O(n)，总计O(nlogn)
\end{itemize}

\begin{lstlisting}[language=C++, style=cppstyle]
int cross(Point p0, Point p1, Point p2) {
    return (p1.x - p0.x) * (p2.y - p0.y) -
           (p1.y - p0.y) * (p2.x - p0.x);}
int dist2(Point p1, Point p2) {
    return (p1.x - p2.x) * (p1.x - p2.x) +
           (p1.y - p2.y) * (p1.y - p2.y);}
bool compare(Point p1, Point p2) {
    int c = cross(p0, p1, p2);
    if (c == 0) return dist2(p0, p1) < dist2(p0, p2);
    return c > 0;}
vector<Point> grahamScan(vector<Point>& points) {
    int n = points.size();
    int p0_idx = 0;
    for (int i = 1; i < n; i++) {
        if (points[i].y < points[p0_idx].y ||
            (points[i].y == points[p0_idx].y &&
             points[i].x < points[p0_idx].x))
            p0_idx = i;}
    swap(points[0], points[p0_idx]);
    p0 = points[0];
    sort(points.begin() + 1, points.end(), compare);
    vector<Point> S;
    S.push_back(points[0]);
    S.push_back(points[1]);
    for (int i = 2; i < n; i++) {
        while (S.size() >= 2 &&
               cross(S[S.size()-2], S.back(), points[i]) <= 0) {
            S.pop_back();}
        S.push_back(points[i]);}
    return S;}
\end{lstlisting}

\begin{enumerate}
    \item 扫描点集找到y最小的点作为p0（有多个时取最左边的），把它交换到数组首位
    \item 从第二个点开始，按"相对p0的极角"排序：用叉积判断，共线时离p0近的排前面
    \item 新建一个栈，先把p0和排序后的第一个点入栈
    \item 从第三个点开始，检查栈顶两个点加上当前点是否"左转"：若不是（右转或共线）就弹出栈顶，重复检查直到满足左转或栈只剩一个点
    \item 把当前点压入栈，继续处理下一个点
    \item 全部处理完后，栈里剩下的点按顺序连起来就是凸包
\end{enumerate}

\par\columnbreak
\section*{2. 两个有序数组的中位数}

设X[0:n-1]和Y[0:n-1]为两个长度均为n的有序数组，设计O(log n)时间的分治算法，找出X和Y的2n个元素的中位数（偶数个元素取中间两个的平均值，奇数个取最中间元素）。

\begin{itemize}
    \item \textbf{问题转化}：2n个元素的中位数是第n小和第n+1小元素的平均值，需同时找到这两个元素
    \item \textbf{分治策略}：对两个数组同时划分，设X的划分边界为i（前i个元素），Y的划分边界为j（前j个元素），要求i+j=n（确保左半部分总长度为n）
    \item \textbf{终止条件}：若X[i-1]$\le$Y[j]且Y[j-1]$\le$X[i]，则左半最大值和右半最小值即为第n小和第n+1小元素
    \item \textbf{范围缩小}：若X[i-1]>Y[j]，说明X的前i个元素过大，需减小i；若Y[j-1]>X[i]，说明Y的前j个元素过大，需减小j
    \item \textbf{时间复杂度}：每次排除一半元素，O(log n)
\end{itemize}

\begin{lstlisting}[language=C++, style=cppstyle]
double findMedianSortedArrays(vector<int>& X, vector<int>& Y) {
    int n = X.size();
    int low = 0, high = n;
    while (low <= high) {
        int i = (low + high) / 2;
        int j = n - i;
        int X_left = (i == 0) ? INT_MIN : X[i - 1];
        int X_right = (i == n) ? INT_MAX : X[i];
        int Y_left = (j == 0) ? INT_MIN : Y[j - 1];
        int Y_right = (j == n) ? INT_MAX : Y[j];
        if (X_left <= Y_right && Y_left <= X_right) {
            if ((n + n) % 2 == 0) {
                return (max(X_left, Y_left) + min(X_right, Y_right)) / 2.0;
            } else {
                return min(X_right, Y_right);}}
        } else if (X_left > Y_right) {
            high = i - 1;
        } else {
            low = i + 1;}}
    return 0.0;}
\end{lstlisting}

\begin{enumerate}
    \item 在较短的数组上二分划分边界i（等长时任选一个），计算另一数组的划分边界j=n-i
    \item 检查边界条件：若X左半最大$\le$Y右半最小且Y左半最大$\le$X右半最小，说明划分正确
    \item 若X左半最大>Y右半最小，说明X的前i个元素太大，让i减小（向左移）
    \item 若Y左半最大>X右半最小，说明Y的前j个元素太大，让i增大（向右移，j随之减小）
    \item 找到正确划分后，偶数个元素取左半最大和右半最小的平均值，奇数个取右半最小
\end{enumerate}

\par\columnbreak
\section*{3. 最大子数组问题}

给定包含n个整数（有正有负）的数组A，设计O(nlogn)时间的分治算法找出和最大的非空连续子数组，同时设计O(n)时间的Kadane算法。

\begin{itemize}
    \item \textbf{分治策略}：将数组A[low..high]拆分为左半A[low..mid]和右半A[mid+1..high]，最大子数组必在左、右、或跨越中点三者之一
    \item \textbf{跨越中点求解}：从mid向左扫描得最大左段和，从mid+1向右扫描得最大右段和，两者相加即为跨越中点的最大子数组和
    \item \textbf{递归求解}：左半和右半的最大子数组通过递归求解，终止条件为子数组长度为1
    \item \textbf{Kadane算法}：维护current\_max（以当前元素结尾的最大子数组和）和global\_max（全局最大），current\_max = max(A[i], current\_max + A[i])
    \item \textbf{时间复杂度}：分治法O(nlogn)，Kadane算法O(n)
\end{itemize}

\begin{lstlisting}[language=C++, style=cppstyle]
int maxCrossingSubarray(vector<int>& A, int low, int mid, int high) {
    int left_sum = INT_MIN, sum = 0;
    for (int i = mid; i >= low; i--) {
        sum += A[i];
        left_sum = max(left_sum, sum);}
    int right_sum = INT_MIN;
    sum = 0;
    for (int i = mid + 1; i <= high; i++) {
        sum += A[i];
        right_sum = max(right_sum, sum);}
    return left_sum + right_sum;}
int maxSubarrayDC(vector<int>& A, int low, int high) {
    if (low == high) return A[low];
    int mid = (low + high) / 2;
    int left_max = maxSubarrayDC(A, low, mid);
    int right_max = maxSubarrayDC(A, mid + 1, high);
    int cross_max = maxCrossingSubarray(A, low, mid, high);
    return max({left_max, right_max, cross_max});}
int maxSubarrayKadane(vector<int>& A) {
    int current_max = A[0];
    int global_max = A[0];
    for (int i = 1; i < A.size(); i++) {
        current_max = max(A[i], current_max + A[i]);
        global_max = max(global_max, current_max);}
    return global_max;}
\end{lstlisting}

\begin{enumerate}
    \item 把数组从中间分开，最大子数组要么全在左边，要么全在右边，要么跨过中点
    \item 左边和右边的最大子数组：递归调用自己去找
    \item 跨中点的最大子数组：从中点往左扫找到最大左段，往右扫找到最大右段，加起来
    \item 三种情况的最大值就是答案
    \item Kadane算法更简单：遍历数组，对于每个位置决定"从这里重新开始"还是"延续之前的子数组"
\end{enumerate}

\par\columnbreak
\section*{4. 循环移位问题}

给定一个按从小到大排序的数组，现将其循环右移若干位（如原数组[1,2,3,4,5]，循环右移2位后为[4,5,1,2,3]），设计算法计算循环右移的位数k（$k \ge 0$），要求时间复杂度尽可能优化。

\begin{itemize}
    \item \textbf{核心性质}：循环右移后的数组是"两段有序子数组"的拼接，第一段所有元素$\ge$第二段所有元素，循环右移位数k等于最小值元素的索引
    \item \textbf{分治查找}：类似折半搜索，比较A[mid]与A[high]，若A[mid]$\le$A[high]则右半有序，最小值在左半；否则最小值在右半
    \item \textbf{终止条件}：low == high时找到最小值，或当A[mid]小于左右邻居时mid即为最小值位置
    \item \textbf{边界处理}：数组未移位（仍完全有序）时，最小值索引为0，k=0
    \item \textbf{时间复杂度}：每次排除一半，O(log n)
\end{itemize}

\begin{lstlisting}[language=C++, style=cppstyle]
int findRotationCount(vector<int>& A) {
    int n = A.size();
    int low = 0, high = n - 1;
    if (A[low] <= A[high]) return 0;
    while (low <= high) {
        int mid = (low + high) / 2;
        int prev = (mid - 1 + n) % n;
        int next = (mid + 1) % n;
        if (A[mid] <= A[prev] && A[mid] <= A[next]) {
            return mid;}
        if (A[mid] <= A[high]) {
            high = mid - 1;
        } else {
            low = mid + 1;}}
    return 0;}
\end{lstlisting}

\begin{enumerate}
    \item 先判断数组是否根本没移位（首元素$\le$尾元素），是的话直接返回0
    \item 找到中点，看它是不是最小值：比左边和右边都小，那就是了
    \item 如果不是，判断哪半边是有序的：右半有序说明最小值在左半，左半有序说明最小值在右半
    \item 向无序的那半边继续找（用二分的方式）
    \item 找到最小值的位置就是循环右移的位数k
\end{enumerate}

\par\columnbreak
\section*{5. 归并排序}

给定数组A，使用归并排序算法对其进行升序排序，要求展示分治的"分解$\to$求解$\to$合并"完整过程。

\begin{itemize}
    \item \textbf{分解}：将数组按索引二分拆分为两个规模相等的子数组，重复拆分直至子数组长度为1
    \item \textbf{基准情况}：子数组长度为1时本身已有序，直接作为子问题的解
    \item \textbf{合并}：对两个有序子数组用双指针法遍历，依次选取较小元素放入临时数组，合并为一个有序数组
    \item \textbf{递归结构}：先递归排序左半，再递归排序右半，最后合并左右
    \item \textbf{时间复杂度}：分解O(1)，合并O(n)，递归树深度logn，总计O(nlogn)
\end{itemize}

\begin{lstlisting}[language=C++, style=cppstyle]
void merge(vector<int>& A, int low, int mid, int high) {
    vector<int> L(A.begin() + low, A.begin() + mid + 1);
    vector<int> R(A.begin() + mid + 1, A.begin() + high + 1);
    int i = 0, j = 0, k = low;
    while (i < L.size() && j < R.size()) {
        if (L[i] <= R[j]) {
            A[k++] = L[i++];
        } else {
            A[k++] = R[j++];}}
    while (i < L.size()) A[k++] = L[i++];
    while (j < R.size()) A[k++] = R[j++];}
void mergeSort(vector<int>& A, int low, int high) {
    if (low < high) {
        int mid = (low + high) / 2;
        mergeSort(A, low, mid);
        mergeSort(A, mid + 1, high);
        merge(A, low, mid, high);}}
\end{lstlisting}

\begin{enumerate}
    \item 把数组从中间分成两半
    \item 对左半部分递归调用归并排序
    \item 对右半部分递归调用归并排序
    \item 合并左右两个有序部分：用两个指针分别指向左右部分的开头，每次选较小的放入原数组
    \item 如果左边或右边还有剩余元素，直接全部复制过去
    \item 重复上述步骤直到数组完全有序
\end{enumerate}

\par\columnbreak
\section*{6. 折半搜索（查找成功案例）}
给定有序数组A，使用折半搜索算法查找目标元素target，要求展示每次查找的low、mid、high值及比较过程，若查找成功返回元素下标，失败返回-1。
\begin{itemize}
    \item \textbf{初始化}：low=0（数组起始），high=n-1（数组结束）
    \item \textbf{中点计算}：mid=$\lfloor$(low+high)/2$\rfloor$
    \item \textbf{比较调整}：若target==A[mid]则查找成功；若target>A[mid]则目标在右半，low=mid+1；若target<A[mid]则目标在左半，high=mid-1
    \item \textbf{终止条件}：找到目标返回mid，或low>high时返回-1
    \item \textbf{时间复杂度}：每次排除一半，O(log n)
\end{itemize}
\begin{lstlisting}[language=C++, style=cppstyle]
int binarySearch(vector<int>& A, int target) {
    int low = 0, high = A.size() - 1;
    while (low <= high) {
        int mid = (low + high) / 2;
        if (A[mid] == target) {
            return mid;
        } else if (A[mid] < target) {
            low = mid + 1;
        } else {
            high = mid - 1;}}
    return -1;}
\end{lstlisting}

\begin{enumerate}
    \item 设置初始搜索范围：low指向第一个元素，high指向最后一个元素
    \item 计算中点位置mid，比较目标值和A[mid]的大小
    \item 如果相等，找到了，返回mid
    \item 如果目标值更大，说明目标在右半边，让low移到mid+1
    \item 如果目标值更小，说明目标在左半边，让high移到mid-1
    \item 重复直到找到目标或low超过high
\end{enumerate}

\par\columnbreak
\section*{7. 折半搜索（查找失败案例）}
给定有序数组A，使用折半搜索算法查找不存在的目标元素，重点关注查找失败时的边界调整及终止条件判断。
\begin{itemize}
    \item \textbf{与成功案例的区别}：算法流程完全一致，区别在于最终触发low>high的终止条件
    \item \textbf{关键观察}：每次比较后缩小搜索范围，当low>high时说明搜索范围为空，目标不存在
    \item \textbf{边界动态}：low和high逐步靠近，最终low=high+1，表示没有元素可查
    \item \textbf{返回值}：失败时返回-1（或其他约定值表示未找到）
    \item \textbf{时间复杂度}：最坏情况需完全排除所有元素，O(log n)
\end{itemize}
\begin{lstlisting}[language=C++, style=cppstyle]
int binarySearchRecursive(vector<int>& A, int target, int low, int high) {
    if (low > high) {
        return -1;}}
    int mid = (low + high) / 2;
    if (A[mid] == target) {
        return mid;
    } else if (A[mid] < target) {
        return binarySearchRecursive(A, target, mid + 1, high);
    } else {
        return binarySearchRecursive(A, target, low, mid - 1);}}}
int binarySearch(vector<int>& A, int target) {
    int low = 0, high = A.size() - 1;
    while (low <= high) {
        int mid = low + (high - low) / 2;
        if (A[mid] == target) return mid;
        else if (A[mid] < target) low = mid + 1;
        else high = mid - 1;}}
    return -1;}
\end{lstlisting}

\begin{enumerate}
    \item 初始化low和high，确定搜索范围
    \item 每次取中点mid，比较目标值和A[mid]
    \item 根据比较结果调整low或high，缩小搜索范围
    \item 当low>high时，说明已经把所有可能的位置都查过了，没找到
    \item 返回-1表示查找失败
\end{enumerate}

\par\columnbreak
\section*{8. 大整数乘法}
设X和Y为两个n位大整数（n为偶数），X拆分为A（前n/2位）和B（后n/2位），Y拆分为C（前n/2位）和D（后n/2位），即$X=A \cdot 2^{n/2}+B$，$Y=C \cdot 2^{n/2}+D$。设计优化方案通过减少乘法次数降低时间复杂度。
\begin{itemize}
    \item \textbf{直接展开}：需计算AC、AD、BC、BD共4次n/2位乘法，递归方程$T(n)=4T(n/2)+O(n)$，复杂度$O(n^2)$未优化
    \item \textbf{Karatsuba算法}：通过代数变换设$Z_1=(A+B)(C+D)=AC+AD+BC+BD$，则$XY=Z_1 \cdot 2^{n/2}+Z_2 \cdot 2^{n/2}+BD$，乘法次数降为3次
    \item \textbf{三次乘法}：只需计算$(A+B)(C+D)$、$AC$、$BD$，通过减法得到$AD+BC=Z_1-AC-BD$
    \item \textbf{递归方程}：$T(n)=3T(n/2)+O(n)$，加法合并时间O(n)
    \item \textbf{时间复杂度}：由Master定理，$a=3$，$b=2$，$\log_b a=\log_2 3 \approx 1.585>1$，故$T(n)=O(n^{\log_2 3}) \approx O(n^{1.585})$
\end{itemize}
\begin{lstlisting}[language=C++, style=cppstyle]
string karatsuba(string X, string Y) {
    int n = max(X.length(), Y.length());
    while (X.length() < n) X = "0" + X;
    while (Y.length() < n) Y = "0" + Y;
    if (n % 2 == 1) {
        n++;
        X = "0" + X;
        Y = "0" + Y;
    }}
    if (n == 1) {
        int val = (X[0] - '0') * (Y[0] - '0');
        return to_string(val);
    }}
    int m = n / 2;
    string A = X.substr(0, m);
    string B = X.substr(m);
    string C = Y.substr(0, m);
    string D = Y.substr(m);
    string Z1 = karatsuba(add(A, B), add(C, D));
    string Z2 = karatsuba(A, C);
    string Z3 = karatsuba(B, D);
    string Z4 = subtract(subtract(Z1, Z2), Z3);
    string result = add(shiftLeft(Z2, n), shiftLeft(Z4, m));
    result = add(result, Z3);
    return result;}
\end{lstlisting}

\begin{enumerate}
    \item 把两个大整数从中间对半拆，得到A、B、C、D四个部分
    \item 直接展开需要AC、AD、BC、BD四次乘法，效率不高
    \item 用Karatsuba技巧：先算(A+B)(C+D)，再算AC和BD，通过减法得到AD+BC
    \item 把三个结果按位权拼接：AC左移n位，(AD+BC)左移n/2位，BD不用移
    \item 把三部分加起来得到最终结果
    \item 比传统方法少一次乘法，复杂度从$O(n^2)$降到$O(n^{1.585})$
\end{enumerate}

\par\columnbreak
\section*{9. Strassen矩阵乘法}

给定两个n$\times$n矩阵A和B，使用Strassen算法（分治法）计算矩阵乘积C=A$\times$B，要求展示7个中间矩阵的定义及复杂度分析。

\begin{itemize}
    \item \textbf{传统方法}：按定义计算C的每个元素需n次乘法，共$n^3$次乘法，递归方程$T(n)=8T(n/2)+O(n^2)$
    \item \textbf{Strassen优化}：构造7个中间矩阵M1~M7，通过7次乘法替代传统的8次乘法
    \item \textbf{7个中间矩阵}：
    \begin{itemize}
        \item $M_1=(A_{11}+A_{22})(B_{11}+B_{22})$，$M_2=(A_{21}+A_{22})B_{11}$
        \item $M_3=A_{11}(B_{12}-B_{22})$，$M_4=A_{22}(B_{21}-B_{11})$
        \item $M_5=(A_{11}+A_{12})B_{22}$，$M_6=(A_{21}-A_{11})(B_{11}+B_{12})$
        \item $M_7=(A_{12}-A_{22})(B_{21}+B_{22})$
    \end{itemize}
    \item \textbf{合并结果}：$C_{11}=M_1+M_4-M_5+M_7$，$C_{12}=M_3+M_5$，$C_{21}=M_2+M_4$，$C_{22}=M_1-M_2+M_3+M_6$
    \item \textbf{时间复杂度}：$T(n)=7T(n/2)+O(n^2)$，Master定理得$O(n^{\log_2 7}) \approx O(n^{2.807})$
\end{itemize}

\begin{lstlisting}[language=C++, style=cppstyle]
vector<vector<int>> matrixAdd(vector<vector<int>>& A,
                              vector<vector<int>>& B) {
    int n = A.size();
    vector<vector<int>> C(n, vector<int>(n));
    for (int i = 0; i < n; i++)
        for (int j = 0; j < n; j++)
            C[i][j] = A[i][j] + B[i][j];
    return C;}
vector<vector<int>> matrixSub(vector<vector<int>>& A,
                              vector<vector<int>>& B) {
    int n = A.size();
    vector<vector<int>> C(n, vector<int>(n));
    for (int i = 0; i < n; i++)
        for (int j = 0; j < n; j++)
            C[i][j] = A[i][j] - B[i][j];
    return C;}
vector<vector<int>> strassen(vector<vector<int>>& A,
                             vector<vector<int>>& B) {
    int n = A.size();
    if (n <= 64) {
        vector<vector<int>> C(n, vector<int>(n, 0));
        for (int i = 0; i < n; i++)
            for (int j = 0; j < n; j++)
                for (int k = 0; k < n; k++)
                    C[i][j] += A[i][k] * B[k][j];
        return C;
    }}
    int m = n / 2;
    return C;}
\end{lstlisting}

\begin{enumerate}
    \item 把大矩阵从中间分成4个小矩阵：A11、A12、A21、A22和B的对应4块
    \item 传统方法需要8次小矩阵乘法才能得到C的4个块
    \item Strassen方法计算7个特殊的中间矩阵M1到M7，每个都是小矩阵的加减后再乘
    \item 通过这7个中间矩阵的加减运算，拼出C的4个块
    \item 虽然加减次数增多，但乘法从8次减到7次，递归下去效率提升明显
    \item 复杂度从$O(n^3)$降到$O(n^{2.807})$，矩阵大时优势显著
\end{enumerate}

\par\columnbreak
\section*{10. 快速排序}
给定数组A，使用快速排序算法对其进行升序排序，要求展示Partition过程及完整的递归排序步骤。
\begin{itemize}
    \item \textbf{Partition过程}：选择基准元素pivot（常选首元素），维护指针p指向小于pivot的区域的末尾，遍历数组使小于pivot的元素交换到前面
    \item \textbf{分区逻辑}：i从1到n-1遍历，若A[i]<pivot则p++，交换A[p]与A[i]，最后交换A[0]与A[p]使pivot归位
    \item \textbf{递归排序}：对pivot左侧子数组A[0..p-1]和右侧子数组A[p+1..n-1]分别递归调用快排
    \item \textbf{终止条件}：子数组长度为0或1时已有序，直接返回
    \item \textbf{时间复杂度}：平均$O(n\log n)$（每次平衡划分），最坏$O(n^2)$（每次划分极不平衡）
\end{itemize}
\begin{lstlisting}[language=C++, style=cppstyle]
int partition(vector<int>& A, int low, int high) {
    int pivot = A[low];
    int p = low;
    for (int i = low + 1; i <= high; i++) {
        if (A[i] < pivot) {
            p++;
            swap(A[p], A[i]);}}
    swap(A[low], A[p]);
    return p;}
void quickSort(vector<int>& A, int low, int high) {
    if (low < high) {
        int p = partition(A, low, high);
        quickSort(A, low, p - 1);
        quickSort(A, p + 1, high);}}
\end{lstlisting}
\begin{enumerate}
    \item 选第一个元素作为基准，用指针p记录"小于基准的区域"的右边界
    \item 从第二个元素开始遍历，如果比基准小，就扩展p，把这个元素交换到前面的区域
    \item 遍历完后，把基准元素交换到p位置，这样基准左边都小，右边都大（或相等）
    \item 对基准左边的子数组递归调用快排
    \item 对基准右边的子数组递归调用快排
    \item 重复直到所有子数组长度不超过1，排序完成
\end{enumerate}

\par\columnbreak
\section*{11. 第k小元素（选择问题）}

给定无序数组A，使用分治法查找该数组的第k小元素，要求展示Partition过程及递归查找步骤。

\begin{itemize}
    \item \textbf{Partition分区}：选择基准元素pivot，划分数组为左（小于pivot）、中（pivot）、右（大于等于pivot）三部分，返回pivot的最终位置pos
    \item \textbf{k值判断}：若k==pos+1则pivot即为第k小元素；若k<pos+1则在左子数组递归查找第k小；若k>pos+1则在右子数组递归查找第k-pos-1小
    \item \textbf{k值调整}：递归右子数组时，k值需减去pos+1（排除左子数组和pivot）
    \item \textbf{平均复杂度}：每次Partition期望划分平衡，$T(n)=T(n/2)+O(n)$，得$O(n)$；最坏$O(n^2)$（每次划分极不平衡）
    \item \textbf{优化}：随机选择pivot或使用Median-of-Medians算法可保证最坏$O(n)$
\end{itemize}

\begin{lstlisting}[language=C++, style=cppstyle]
int partition(vector<int>& A, int low, int high) {
    int pivot = A[low];
    int p = low;
    for (int i = low + 1; i <= high; i++) {
        if (A[i] < pivot) {
            p++;
            swap(A[p], A[i]);}}
    swap(A[low], A[p]);
    return p;}
int quickSelect(vector<int>& A, int low, int high, int k) {
    if (low == high) return A[low];
    int pos = partition(A, low, high);
    if (k == pos + 1) {
        return A[pos];
    } else if (k < pos + 1) {
        return quickSelect(A, low, pos - 1, k);
    } else {
        return quickSelect(A, pos + 1, high, k - pos - 1);}}}
int findKthSmallest(vector<int>& A, int k) {
    return quickSelect(A, 0, A.size() - 1, k);}
\end{lstlisting}

\begin{enumerate}
    \item 选一个基准元素，用Partition把数组分成"小于基准"和"大于等于基准"两部分
    \item 基准元素归位后得到位置pos，说明它左边有pos个元素比它小
    \item 如果k等于pos+1，那基准就是第k小元素，直接返回
    \item 如果k小于pos+1，说明第k小元素在左边，递归在左边找第k小
    \item 如果k大于pos+1，说明第k小元素在右边，递归在右边找第k-pos-1小（要减去左边和基准）
    \item 重复直到找到为止
\end{enumerate}

\par\columnbreak
\section*{12. 堆排序}

给定数组A，使用堆排序算法对其进行升序排序，要求展示最大堆的建立过程、堆顶元素的弹出与堆调整过程。

\begin{itemize}
    \item \textbf{堆性质}：最大堆中每个节点值$\ge$其子节点值，数组表示中节点i的左子为2i+1，右子为2i+2，父为$\lfloor$(i-1)/2$\rfloor$
    \item \textbf{建堆}：从最后一个非叶子节点$\lfloor$n/2$\rfloor$-1开始，向前依次执行Heapify向下调整，时间$O(n)$
    \item \textbf{Heapify}：对节点i，若其值小于子节点，则与最大子节点交换，并递归调整被交换的子节点
    \item \textbf{排序循环}：交换堆顶（最大值）与堆尾，缩小堆范围，对新的堆顶执行Heapify重建最大堆
    \item \textbf{时间复杂度}：建堆$O(n)$，n次弹出各$O(\log n)$，总计$O(n\log n)$；空间$O(1)$原地排序
\end{itemize}

\begin{lstlisting}[language=C++, style=cppstyle]
void heapify(vector<int>& A, int n, int i) {
    int largest = i;
    int left = 2 * i + 1;
    int right = 2 * i + 2;
    if (left < n && A[left] > A[largest])
        largest = left;
    if (right < n && A[right] > A[largest])
        largest = right;
    if (largest != i) {
        swap(A[i], A[largest]);
        heapify(A, n, largest);}}
void buildHeap(vector<int>& A, int n) {
    for (int i = n / 2 - 1; i >= 0; i--) {
        heapify(A, n, i);}}
void heapSort(vector<int>& A) {
    int n = A.size();
    buildHeap(A, n);
    for (int i = n - 1; i > 0; i--) {
        swap(A[0], A[i]);
        heapify(A, i, 0);}}}
\end{lstlisting}

\begin{enumerate}
    \item 从最后一个非叶子节点开始，对每个节点执行Heapify调整，把数组变成最大堆
    \item Heapify调整时，如果节点比子节点小，就跟最大子节点交换，然后继续往下调整
    \item 建好堆后，堆顶就是最大值，把它和堆尾元素交换，最大值就放到有序区了
    \item 堆范围缩小1，对新的堆顶再执行Heapify重建最大堆
    \item 重复"交换堆顶和堆尾、缩小堆、重建堆"这个过程
    \item 直到堆只剩一个元素，排序完成
\end{enumerate}

\par\columnbreak
\section*{13. 最近点对问题}

给定平面上n个点的集合Q，设计分治算法找出距离最近的两个点，要求说明预处理、分解、递归求解、合并的完整流程。

\begin{itemize}
    \item \textbf{预处理}：将点集按x坐标排序，便于均匀划分，时间$O(n\log n)$
    \item \textbf{分解}：按x坐标中位数用垂线x=m划分点集为L（x$\le$m）和R（x>m），确保规模平衡
    \item \textbf{递归求解}：分别找L和R的最近点对，距离为d1和d2，令d=min(d1,d2)
    \item \textbf{合并（临界区）}：筛选x$\in$[m-d,m+d]的点集Yd按y排序，对每个点仅与后6个点比较（鸽巢原理），时间$O(n)$
    \item \textbf{时间复杂度}：递归方程$T(n)=2T(n/2)+O(n)$，Master定理得$O(n\log n)$
\end{itemize}

\begin{lstlisting}[language=C++, style=cppstyle]
struct Point {
    double x, y;
double dist(Point p1, Point p2) {
    return sqrt((p1.x - p2.x) * (p1.x - p2.x) +
                (p1.y - p2.y) * (p1.y - p2.y));}
double bruteForce(vector<Point>& pts, int left, int right) {
    double minDist = INT_MAX;
    for (int i = left; i <= right; i++)
        for (int j = i + 1; j <= right; j++)
            minDist = min(minDist, dist(pts[i], pts[j]));
    return minDist;}
double stripClosest(vector<Point>& strip, double d) {
    double minDist = d;
    sort(strip.begin(), strip.end(),
         [](Point a, Point b) { return a.y < b.y; });
    for (size_t i = 0; i < strip.size(); i++) {
        for (size_t j = i + 1; j < strip.size() &&
             (strip[j].y - strip[i].y) < minDist; j++) {
            minDist = min(minDist, dist(strip[i], strip[j]));
        }    }}
    return minDist;}
double closestUtil(vector<Point>& pts, int left, int right) {
    if (right - left + 1 <= 3)
        return bruteForce(pts, left, right);
    int mid = (left + right) / 2;
    Point midPoint = pts[mid];
    double d1 = closestUtil(pts, left, mid);
    double d2 = closestUtil(pts, mid + 1, right);
    double d = min(d1, d2);
    vector<Point> strip;
    for (int i = left; i <= right; i++)
        if (abs(pts[i].x - midPoint.x) < d)
            strip.push_back(pts[i]);
    return min(d, stripClosest(strip, d));}}}
double closestPair(vector<Point>& pts) {
    sort(pts.begin(), pts.end(),
         [](Point a, Point b) { return a.x < b.x; });
    return closestUtil(pts, 0, pts.size() - 1);}
\end{lstlisting}

\begin{enumerate}
    \item 先把所有点按x坐标排序，方便后面均匀划分
    \item 从x坐标中位数位置把点集分成左右两半，递归在两边找最近的点对
    \item 得到两边的最小距离d后，最近点对要么在左边，要么在右边，要么跨过分界线
    \item 跨界点对只能在分界线左右各d宽度的条形区域里
    \item 把这个条形区域的点按y坐标排序，每个点只需跟后面6个点比较距离（鸽巢原理保证）
    \item 三种情况（左、右、跨）取最小就是最终答案
\end{enumerate}

\par\columnbreak
\section*{14. 天际轮廓问题}

给定n座建筑物，每个建筑物用三元组(Li,Ri,Hi)表示（Li为左下x坐标，Ri为右下x坐标，Hi为高度），设计O(nlogn)算法求出这n座建筑物的天际轮廓。

\begin{itemize}
    \item \textbf{天际轮廓表示}：用高度变化的关键点序列表示，如单个建筑物的轮廓为[(Li,0),(Li,Hi),(Ri,Hi),(Ri,0)]
    \item \textbf{分治分解}：将建筑物集均分为左右两个子集，递归求得左右
    \item \textbf{双指针合并}：用双指针i、j分别遍历左右天际线，按x坐标升序合并，同时维护当前左右高度hL、hR，取max(hL,hR)作为合并后当前高度
    \item \textbf{关键点记录}：当当前max(hL,hR)与上次记录的高度不同时，记录新的关键点(x, max(hL,hR))
    \item \textbf{时间复杂度}$T(n)=2T(n/2)+O(n)$，得$O(n\log n)$
\end{itemize}

\begin{lstlisting}[language=C++, style=cppstyle]
struct Point {
    int x, height;
    bool isStart;
};
vector<Point> mergeSkylines(vector<Point>& left, vector<Point>& right) {
    vector<Point> result;
    int i = 0, j = 0;
    int hL = 0, hR = 0, currH = 0;
    while (i < left.size() && j < right.size()) {
        int x;
        if (left[i].x < right[j].x ||
            (left[i].x == right[j].x && left[i].height >= right[j].height)) {
            hL = left[i].height;
            x = left[i].x;
            i++;
        } else {
            hR = right[j].height;
            x = right[j].x;
            j++;}}
        int newH = max(hL, hR);
        if (result.empty() || newH != currH) {
            result.push_back({x, newH});
            currH = newH;}}
    while (i < left.size()) {
        if (left[i].height != currH) {
            result.push_back(left[i]);
            currH = left[i].height;}}
        i++;}
    while (j < right.size()) {
        if (right[j].height != currH) {
            result.push_back(right[j]);
            currH = right[j].height;}}
        j++;}
    return result;}
vector<Point> getSkyline(vector<vector<int>>& buildings, int low, int high) {
    if (low == high) {
        return {{buildings[low][0], buildings[low][2]},
                {buildings[low][1], 0}};}}
    int mid = (low + high) / 2;
    auto left = getSkyline(buildings, low, mid);
    auto right = getSkyline(buildings, mid + 1, high);
    return mergeSkylines(left, right);}}
\end{lstlisting}

\begin{enumerate}
    \item 如果只有一个建筑，天际线就是它的4个顶点
    \item 多个建筑时，从中间分成左右两半，递归求各自的天际线
    \item 用双指针遍历两条天际线，x坐标小的先处理
    \item 维护当前左右天际线的高度，取最大值为合并后当前高度
    \item 只要当前最大高度和上次记录不一样，就记录新的关键点
    \item 重复直到处理完所有关键点，最终得到的就是完整天际线
\end{enumerate}

\par\columnbreak
\section*{15. 两元素和为X的分治算法}

给定由n个实数构成的集合S和实数X，判断S中是否存在两个元素的和为X，要求设计分治算法求解。

\begin{itemize}
    \item \textbf{预处理}：先用归并排序对S排序，时间$O(n\log n)$，为后续双指针查找做准备
    \item \textbf{分治分解}：将有序集合S均分为左子集L和右子集R
    \item \textbf{递归求解}：递归判断L和R内部是否存在两元素和为X
    \item \textbf{合并验证}：用双指针i（指向L首）、j（指向R尾）遍历，判断是否存在跨子集的两元素和为X
    \item \textbf{双指针逻辑}：若S[i]+S[j]==X则返回true；若<X则i++（增大和）；若>X则j--（减小和）
    \item \textbf{时间复杂度}：排序$O(n\log n)$，递归方程$T(n)=2T(n/2)+O(n)$，总计$O(n\log n)$
\end{itemize}

\begin{lstlisting}[language=C++, style=cppstyle]
bool hasPairCross(vector<int>& S, int low, int mid, int high, int X) {
    int i = low;
    int j = high;
    while (i <= mid && j > mid) {
        int sum = S[i] + S[j];
        if (sum == X) {
            return true;
        } else if (sum < X) {
            i++;
        } else {
            j--;}}}}
    return false;}
bool hasTwoSumDC(vector<int>& S, int low, int high, int X) {
    if (low >= high) return false;
    int mid = (low + high) / 2;
    bool leftFound = hasTwoSumDC(S, low, mid, X);
    bool rightFound = hasTwoSumDC(S, mid + 1, high, X);
    if (leftFound || rightFound) return true;
    return hasPairCross(S, low, mid, high, X);}}
bool hasTwoSum(vector<int>& S, int X) {
    sort(S.begin(), S.end());
    return hasTwoSumDC(S, 0, S.size() - 1, X);}}}
\end{lstlisting}

\begin{enumerate}
    \item 先把数组排序（用归并排序），为双指针查找做准备
    \item 把数组从中间分成左右两半，递归判断两边内部有没有符合条件的
    \item 对于跨两边的情况，用双指针：一个指左半开头，一个指右半结尾
    \item 计算两个指针指向的元素的和，如果等于X就找到了
    \item 如果和小于X，左指针右移（找个更大的）；如果大于X，右指针左移（找个更小的）
    \item 三种情况（左内、右内、跨）只要有一个true就返回true
\end{enumerate}

\par\columnbreak
\section*{16. 逆序对的分治算法}

给定实数序列A，若i<j且A[i]>A[j]，则(i,j)为逆序对。要求用分治方法求序列中逆序对的个数。

\begin{itemize}
    \item \textbf{逆序对定义}：前大后小的元素对，如[3,1,2]的逆序对为(3,1)、(3,2)
    \item \textbf{分治分解}：将序列均分为左子序列L和右子序列R
    \item \textbf{递归求解}：递归计算L和R内部的逆序对个数countL和countR
    \item \textbf{合并统计}：归并L和R为有序序列，同时统计跨L、R的逆序对个数cross
    \item \textbf{跨区间统计}：归并时若L[i]>R[j]，则L[i]及L[i]之后的所有元素都与R[j]构成逆序对，count+=mid-i+1
    \item \textbf{时间复杂度}：归并排序的合并步骤嵌入统计，$T(n)=2T(n/2)+O(n)$，总计$O(n\log n)$
\end{itemize}

\begin{lstlisting}[language=C++, style=cppstyle]
int mergeAndCount(vector<int>& A, int low, int mid, int high) {
    vector<int> L(A.begin() + low, A.begin() + mid + 1);
    vector<int> R(A.begin() + mid + 1, A.begin() + high + 1);
    int i = 0, j = 0, k = low;
    int count = 0;
    while (i < L.size() && j < R.size()) {
        if (L[i] <= R[j]) {
            A[k++] = L[i++];
        } else {
            A[k++] = R[j++];
            count += L.size() - i;
        }    }}
    while (i < L.size()) A[k++] = L[i++];
    while (j < R.size()) A[k++] = R[j++];
    return count;}
int countInversions(vector<int>& A, int low, int high) {
    if (low >= high) return 0;
    int mid = (low + high) / 2;
    int countL = countInversions(A, low, mid);
    int countR = countInversions(A, mid + 1, high);
    int countCross = mergeAndCount(A, low, mid, high);
    return countL + countR + countCross;}}}
int countInversions(vector<int>& A) {
    return countInversions(A, 0, A.size() - 1);}}}
\end{lstlisting}

\begin{enumerate}
    \item 把序列从中间分成左右两半
    \item 递归统计左半内部的逆序对个数
    \item 递归统计右半内部的逆序对个数
    \item 归并左右两半为有序序列，同时统计跨区间的逆序对
    \item 统计跨区间时，如果左边元素大于右边，那左边剩下的元素也都大于右边当前元素
    \item 总逆序对数等于"左内+右内+跨区间"之和
\end{enumerate}

\par\columnbreak
\section*{17. 数组中两元素和大于0的对数量}

给定大小为n的数组A，求数组中不同对(i,j)（i<j）的数量，使得A[i]+A[j]>0。

\begin{itemize}
    \item \textbf{分治排序}：用归并排序将A升序排列，时间$O(n\log n)$，为双指针统计做准备
    \item \textbf{分治分解}：将有序数组划分为左子集L和右子集R
    \item \textbf{递归统计}：分别计算L内部、R内部符合条件的对数量countL和countR
    \item \textbf{合并统计}：用双指针i（指向L首）、j（指向R尾）遍历，统计跨子集的对数量cross
    \item \textbf{双指针逻辑}：若A[i]+A[j]>0，则R[j]及R[j]之前的所有元素与A[i]的和都>0，count+=j-mid+1，i++；否则j--
    \item \textbf{时间复杂度}：排序$O(n\log n)$，分治递归方程$T(n)=2T(n/2)+O(n)$，总计$O(n\log n)$
\end{itemize}

\begin{lstlisting}[language=C++, style=cppstyle]
int countCross(vector<int>& A, int low, int mid, int high) {
    int i = low;
    int j = high;
    int count = 0;
    while (i <= mid && j > mid) {
        if (A[i] + A[j] > 0) {
            count += j - mid;
            i++;
        } else {
            j--;}}}}
    return count;}
int countPairsDC(vector<int>& A, int low, int high) {
    if (low >= high) return 0;
    int mid = (low + high) / 2;
    int countL = countPairsDC(A, low, mid);
    int countR = countPairsDC(A, mid + 1, high);
    int countCross = countCross(A, low, mid, high);
    return countL + countR + countCross;}}}
int countPairsWithSumGreaterThanZero(vector<int>& A) {
    sort(A.begin(), A.end());
    return countPairsDC(A, 0, A.size() - 1);}}}
\end{lstlisting}

\begin{enumerate}
    \item 先用归并排序把数组升序排列
    \item 把数组从中间分成左右两半，递归统计两边内部符合条件的对数
    \item 对于跨两边的情况，用双指针：一个指左半开头，一个指右半结尾
    \item 如果两个指针指向的元素和大于0，那右半剩下的所有元素跟左半当前元素的和也都大于0
    \item 累加符合条件的对数，然后左指针右移（尝试更小的元素）
    \item 如果和小于等于0，右指针左移（尝试更大的元素）
    \item 总对数等于"左内+右内+跨子集"之和
\end{enumerate}

\par\columnbreak

% ========== 分治算法方法论总结 ==========
{\large\textbf{分治算法解题的"固有方法论"}}
\vspace{0.2cm}

分治算法（Divide and Conquer）的核心思想是"\textbf{大事化小，小事化了，最后合并}"。在解决此类题目时，可以遵循标准的\textbf{"三步走"}策略：

\textbf{1. 分解 (Divide)}：将原问题划分成若干个规模较小、相互独立、与原问题形式相同的子问题。\\
\textcolor{gray}{$\bullet$ 代码体现：通常是寻找中点 \texttt{mid}，将数组下标范围划分为 \texttt{[left, mid]} 和 \texttt{[mid+1, right]}。}

\textbf{2. 解决 (Conquer)}：递归地求解这些子问题。\\
\textcolor{gray}{$\bullet$ 基准情况 (Base Case)：必须定义递归结束的条件（例如数组长度为1或2时），直接求解。}\\
\textcolor{gray}{$\bullet$ 代码体现：递归调用函数本身，如 \texttt{solve(left, mid)}。}

\textbf{3. 合并 (Combine)}：将子问题的解合并，得到原问题的解。这是分治算法中最关键、最难设计的一步。\\
\textcolor{gray}{$\bullet$ 代码体现：在递归返回后，编写逻辑将左右两边的结果整合成一个结果。}

\vspace{0.3cm}
\textbf{复杂度分析套路：}

通常使用主定理（Master Theorem）或递归树。形式为 $T(n) = aT(n/b) + f(n)$。

\textcolor{gray}{$\bullet$ $a$：子问题个数}\\
\textcolor{gray}{$\bullet$ $b$：子问题缩小的倍数}\\
\textcolor{gray}{$\bullet$ $f(n)$：合并步骤的代价}

% ========== 题目内容结束 ==========

\end{multicols*}

\end{document}
